\section{A frozen shared proposal on unseen workloads}\label{sec:sharedstudy89}
The joint-domain construction in Section~\ref{sec:population86} establishes a mathematical shared-operator guarantee. A trained proposal's incremental numerical value is a different question. We test that value in the original investment model without changing its continuous innovation law, feasible action cap or protected Bellman certificate. This experiment adds a prospective comparison to the inherited recursive-query and complete-policy studies; it does not substitute one of those objects for another.

\subsection{Construction, selection and evaluation}
The protocol is fixed before fitting. The training stream contains 128 states in each of dimensions two and eight, with two remaining dates. Protected adaptive search at tolerance $1/256$ supplies feasible action labels, normalized by the original action cap. Both approximation families receive the same fourteen features: eight zero-padded state coordinates, the overall, even-coordinate and odd-coordinate means, the mean square, the mean adjacent product, and the dimension divided by eight. The neural model is one width-32 ReLU network shared across both dimensions. The nonneural comparator is one complete quadratic on the same features. Each family has one attempted fit and fixed optimization settings, with no hyperparameter search. All coefficients, training rows and fitting diagnostics are retained before calibration.

Calibration has 64 new states in each training dimension. The five services are the adaptive baseline, bisection, a one-date analytic proposal, the shared ReLU proposal and the shared quadratic proposal. Each runs at tolerance $1/64$ using the unchanged protected oracle. Selection minimizes their total action-value queries, with the prespecified tie order. A terminal query is already included in that total and is not added a second time. Calibration selects the adaptive baseline: its 990 queries compare with 999 for bisection, 1008 for the analytic proposal, 996 for ReLU and 999 for the quadratic proposal. This choice is committed before the evaluation states are generated.

Evaluation contains 128 fresh states in each of dimensions two, four and eight. Dimension four is absent from both training and calibration. All services receive the identical states; three fresh-process repetitions preserve the complete timing distribution. The saved coefficients and evaluation inputs also define the cross-architecture replay. Repetitions and machine replay do not add independent states. All streams are reproducible seeded numerical workloads. A seed alone is not an assumption establishing population randomness; Proposition~\ref{prop:selection86} and the coverage theorem state the separate assumptions needed for probabilistic extrapolation.

\input{tables/work89}

\subsection{Results and their economic interpretation}
The two fitting attempts return finite coefficients without a recorded warning. All held-out services meet the common certificate tolerance. The adaptive service attains the smallest total query count in each evaluation dimension; at dimension eight it ties bisection and the quadratic proposal. Its empirical query regret against the declared menu is therefore zero in each dimension. This is not a zero expected-regret theorem and not a comparison with the original Bellman optimum.

The ReLU proposal uses 987, 1017 and 987 queries in dimensions two, four and eight, respectively, against 975, 993 and 969 for the adaptive service. The predictor is genuinely shared and is evaluated on unseen states and an unseen dimension, but it does not reduce verification work in this experiment. The protected first partition already resolves many workloads before any model call. In the remaining workloads, a model's proposed query can fail to improve the protected baseline's allocation. All such calls and refinements are charged. This result identifies the distinction between a valid common operator, a successfully returned fit and a cost-saving proposal; none implies the next.

Complete process times include interpreter startup, coefficient loading, computation, serialization and shutdown. The numerical appendix reports every repetition and peak resident memory on each recorded architecture, rather than only the fastest execution. Logical dense-model arithmetic and disjoint oracle queries remain separate accounts. They are not measured CPU instructions, energy expenditures or monetary costs. Since the calibration objective prices only action-value queries, its selected service need not minimize process time or a different priced resource vector. The price-region analysis in Section~\ref{sec:prices86} gives the appropriate conditional decision rule once those prices and complete resource accounts are specified.

\paragraph{Protection is not prediction success.}
The diagnostic catalogue separately records prediction-free returns, attempted screens, screen closures, protected refinements, invalid predictions, feasibility repairs and zero-progress probes. No failed fit or invalid prediction is omitted; the prescribed safe fallback remains a service outcome, not successful neural prediction. The two fits in this experiment do not erase the inherited nonconvergence warnings or estimate their population frequency. Original-law policy-cost conclusions continue to come from the separately frozen nine-controller evaluation in Section~\ref{sec:prospective86}, not from this query-proposal table.
